% LOI TUA --- Quyển XXV V2
\chapter*{Lời Tựa}
\addcontentsline{toc}{chapter}{Lời Tựa}
\markboth{Lời Tựa}{Lời Tựa}

\begin{truyen}{Lời Tựa}{Rawer Classroom Reality}
\chuhoa{B}{ản V2 này không cố làm lớp học nghe đẹp. Nó cố làm lớp học nghe đúng. Đúng cái kiểu học sinh đang bí thì chống chế, đang quê thì cãi, đang mệt thì nói ngang, đang sợ thì nói trớ.}
\textit{(This V2 does not try to make classroom life sound pretty. It tries to make it sound true: students excuse, argue, dodge, or speak sideways when they are embarrassed, tired, or afraid.)}

Quyển cũ giữ tiếng cười. Quyển này giữ thêm độ gai. Vì ngoài đời, nhiều câu nói trong lớp không tròn trịa như sách đạo lý. Nó cụt, nó khựng, nó hơi khó nghe, nhưng lại rất thật và rất hay gặp.
\textit{(The earlier version kept the humor. This one keeps more friction. Real classroom speech is often blunt, awkward, and not polished, yet very familiar.)}

Trong 10 chương và 40 màn đối đáp, cuốn này đi thẳng vào những tình huống học sinh hay dùng lý sự cùn nhất: lười nhưng muốn giữ sĩ diện, sợ điểm thấp nhưng không dám nhận mình hổng, muốn được hiểu nhưng lại chọn cách chống nạnh bằng lời nói.
\textit{(Across 10 chapters and 40 exchanges, this book goes straight into the most common blunt excuses: laziness protected by pride, fear of low grades hidden behind denial, and a wish to be understood expressed in the wrong tone.)}

Nếu bạn là giáo viên, mong bạn thấy rõ hơn cái phần ``gồ ghề'' của học trò để đáp lại cho đúng. Nếu bạn là học sinh, mong bạn đọc xong sẽ nhận ra nhiều câu mình từng nói tưởng khôn mà thật ra chỉ đang che một chỗ yếu chưa dám nhìn.
\textit{(If you are a teacher, may you better see the rough edges behind student behavior. If you are a student, may you notice how some clever-sounding excuses only hide a weakness you have not faced yet.)}
\end{truyen}

\begin{baihoc}
Nói thật về lớp học không làm mất đi sự tử tế. Nó chỉ làm bài học bớt giả.
\textit{(Speaking honestly about school does not remove kindness. It removes falseness.)}
\end{baihoc}

\begin{ghinhoanh}
\textbf{Short English to remember:}

Raw words often hide real pressure.

Honest teaching is not always soft teaching.
\end{ghinhoanh}

\begin{tuhocnhanh}
1. Em hay chống chế vì lười, vì sợ, hay vì sĩ diện? \textit{(Do you excuse yourself out of laziness, fear, or pride?)}

2. Một câu em từng nói để đỡ quê là gì? \textit{(What is one thing you once said just to save face?)}

3. Em muốn người lớn hiểu mình ở điểm nào nhất? \textit{(What do you most want adults to understand about you?)}
\end{tuhocnhanh}

\begin{flushright}
\textbf{Tôi Nguyễn Văn Sang}\\
\textit{GV THPT Nguyễn Hữu Cảnh - TP HCM}
\end{flushright}

\ngancach